\chapter{Análise de Dados}
\label{ch:analise-dados}

O processo de avaliação por parte dos professores é fundamental para validar a viabilidade, a aplicabilidade e a relevância do sistema desenvolvido, garantindo que o artefato pedagógico atenda às necessidades educacionais e seja eficaz na promoção do aprendizado. A sessão foi realizada em ambiente virtual (on-line), formato que viabiliza a comunicação síncrona em tempo real, elimina barreiras geográficas e propicia maior espontaneidade nas interações \cite{abreu2009}. Seguindo o procedimento metodológico para grupos focais virtuais, a atividade iniciou-se com uma demonstração dirigida do sistema, na qual o moderador apresentou os objetivos e o funcionamento da plataforma para alinhar o entendimento dos participantes \cite{abreu2009}.

Para a condução desta etapa, utilizou-se a técnica de grupo focal, que possibilita compreender, de forma aprofundada, as percepções, experiências e significados atribuídos pelos participantes ao fenômeno estudado. Por se tratar de uma abordagem qualitativa, o grupo focal permite captar nuances, contradições e consensos construídos coletivamente, oferecendo uma compreensão interpretativa mais ampla do problema investigado. O planejamento desta técnica é rigoroso, exigindo a seleção de participantes cujos perfis estejam em consonância com os objetivos da pesquisa \cite{gil2021pesquisa}.

A validação contou com a participação de professores de matemática do curso de Sistemas de Informação, cuja expertise técnica e pedagógica qualifica as percepções coletadas. O objetivo foi garantir que as intervenções educacionais ultrapassassem a idealização teórica, verificando a capacidade do sistema de operar efetivamente diante dos desafios e restrições reais do contexto educacional \cite{plomp2018pesquisa}. Para tanto, foram empregados critérios de avaliação fundamentados em indicadores qualitativos, com foco na exatidão dos conteúdos matemáticos (enunciados, dados e gabaritos), na eficiência do processamento das demandas cognitivas e na capacidade das mecânicas de gamificação em estimular o raciocínio.

Após a apresentação, as discussões ocorreram de forma livre e interativa, resultando em percepções relevantes sobre a usabilidade e aplicabilidade do artefato. Paralelamente, foram evidenciados aspectos positivos da solução, como o papel do erro como elemento de reforço no processo de aprendizagem, incentivando o estudante a refazer os passos e analisar criticamente a resolução. Os participantes ressaltaram o papel do aluno na auto-análise e na reflexão perante ao erro, destacando que o estudante deve reconhecer o próprio equívoco e buscar sua correção de forma autônoma. As dicas do sistema também foram mencionadas como suporte relevante para orientar o estudante nesse processo, sem revelar a solução imediata, preservando o esforço cognitivo do aluno. Observou-se que as mecânicas de gamificação implementadas, especificamente a fragmentação da questão em níveis, a penalização de \gls{xp} por erros recorrentes e o custo estratégico associado ao uso excessivo de auxílios, forçam o usuário a um processo de autorreflexão sobre o erro. Os participantes também destacaram a importância de se aplicar rigorosamente os conceitos de abstração, decomposição e reconhecimento de padrões, bem como a necessidade de estruturar a sequência lógica das soluções.

No que tange às melhorias, destacaram-se a necessidade de aprimorar a interação com o sistema, a reestruturação da etapa de abstração para foco nos pontos centrais do problema e o refinamento da decomposição para uma fragmentação mais analítica da problemática. Adicionalmente, sugeriu-se a implementação de mecânicas de passos sequenciais nas fases de decomposição e algoritmo para melhor estruturação do raciocínio lógico do estudante, bem como a criação de um folhetim de uso do sistema, contendo um passo a passo para guiar o usuário durante a interação com a plataforma. Outro ponto de aprimoramento reside na possibilidade de implementar mecanismos de orientação direcionada, que priorizem as fases com maior índice de erro em detrimento daquelas com maior taxa de acerto nas sessões, compelindo o usuário a aprimorar seu desempenho nas etapas de maior dificuldade. Os participantes também apontaram a falta de um retorno mais específico para o aluno, uma vez que o sistema é determinístico e apresenta as mesmas respostas independentemente do contexto do estudante, o que limita a personalização da orientação durante a resolução dos problemas.

\section{Associação do relato dos professores ao referencial teórico}

Os depoimentos dos professores de matemática convergem com o referencial teórico apresentado na Fundamentação Teórica em diversos aspectos. Em primeiro lugar, a valorização do erro como elemento de reforço no processo de aprendizagem, na qual o estudante é incentivado a refazer os passos e analisar criticamente a resolução, ecoa diretamente a posição de \cite{piton2018analise}, para quem a compreensão adequada do enunciado e a superação dos erros são centrais para o desempenho matemático, bem como a abordagem de \cite{wing2006}, que reconhece o erro como parte intrínseca do processo de resolução de problemas.

O destaque dado pelos participantes às mecânicas de gamificação, como fragmentação da questão em níveis, penalização de \gls{xp} por erros recorrentes e custo estratégico do uso excessivo de auxílios, dialoga com o referencial de \cite{silpasuwanchai2016developing}, que aponta pontos, recompensas e desafios como elementos que elevam o engajamento comportamental e cognitivo, e com \cite{michael2006seriousgame}, que enfatiza o retorno crítico constante, a correção rápida de erros e o aprendizado por tentativa e erro como motores do envolvimento do estudante. Ao forçar o usuário a uma autorreflexão sobre o erro, o jogo propiciado reproduz, em ambiente lúdico, o ciclo de tentativa, fracasso e correção que o referencial identifica como condição para a construção do conhecimento.

A importância atribuída pelos professores aos conceitos de abstração, decomposição e reconhecimento de padrões, bem como à estruturação de uma sequência lógica de soluções, encontra respaldo em \cite{wing2006}, que define essas competências como o núcleo do pensamento computacional aplicado à resolução de problemas matemáticos. A sugestão de implementar passos sequenciais nas fases de decomposição e algoritmo reforça a necessidade de guiar o estudante na construção do raciocínio lógico, ponto que também aparece em \cite{pontes2019questoes}, ao defender o estímulo a estratégias corretas de resolução em detrimento da mera repetição de fórmulas.

Por fim, a demanda por mecanismos de orientação direcionada, que priorizam as fases com maior índice de erro, e a criação de um guia de uso (folhetim) correspondem à proposta do referencial de oferecer suporte estruturado ao estudante, alinhando-se ao eixo do Pensamento Computacional da \gls{bncc}, que defende a modelagem e a resolução metódica de problemas com apoio de tecnologias digitais \cite{brasil2017,brasil_bncc_computacao}. Ademais, o relato dos professores sobre o papel do aluno na auto-análise e na reflexão perante ao erro, bem como sobre a importância das dicas do sistema como suporte que orienta sem revelar a solução imediata, dialoga com a concepção de erro como oportunidade de aprendizagem de \cite{piton2018analise} e com o retorno crítico constante e o aprendizado por tentativa e erro defendidos por \cite{michael2006seriousgame}. Dessa forma, o relato dos professores valida a aplicabilidade do artefato e aponta ajustes que mantêm o sistema ancorado nos princípios teóricos que fundamentam a pesquisa.
